\vssub
\subsection{~风和流} \label{sub:num_w_c}
\vssub

\noindent
模式输入主要由风场和流场组成。在模式内部，风和流在每个时间步长 $\Delta t_g$ 更新，并表示所考虑时间步末端的取值。模式提供若干插值方法（编译时选择）。默认情况下，时间插值采用速度和方向的线性插值（使风或流按最小角度转向）。也可选择对风速或流速进行修正，以近似守恒能量而不是守恒风速。相应修正因子 $X_u$ 计算为

% eq:X_u10

\begin{equation}
X_u = \max \left [ \: 1.25 \: , \: \frac{u_{10,rms}}{u_{10,l}}
\right ] \: , \label{eq:X_u10} \end{equation}

\noindent
其中 $u_{10,l}$ 是线性插值得到的速度，$u_{10,rms}$ 是 rms 插值得到的速度。最后，也可选择使风保持常值并以不连续方式改变（该选项不适用于流）。

注意，\ws{} 的辅助程序包括一个用于预处理输入场的程序（见 \para\ref{sec:ww3prep}）。该程序会将网格化场转换到波浪模式网格上。对于风和流，该程序使用矢量分量的双线性插值。利用类似于式~(\ref{eq:X_u10}) 的修正因子，可对这种插值进行修正，以近似守恒风或流的速度或能量。
